% ID: SEL-D0029-OB-E03
% ESTADO: cotejada. Segunda pasada visual completa y borrador compilado revisados.
% Fuente: Selectividad/matematicas_aplicadas/suplente septiembre/Examen-Matemáticas Aplicadas a las Ciencias Sociales.pdf
% SHA-256: 4b8884e307df4f5ddaf079e0f10be3d3f74d4793874eb03b355d34f4856579cb
% Página: 2; opción B; ejercicio original 3.

En una concentración de 250 deportistas hay 120 que juegan al fútbol, 60 que juegan al tenis y 70 que juegan al baloncesto. El 75\% de los que juegan al fútbol, el 65\% de los que juegan al tenis y el 60\% de los que juegan al baloncesto son además aficionados al ciclismo.

Se selecciona al azar uno de los deportistas.
\begin{enumerate}[label=\alph*)]
\item \textbf{(1.5 puntos)} \textquestiondown{}Cuál es la probabilidad de que sea aficionado al ciclismo?
\item \textbf{(1 punto)} Si es aficionado al ciclismo, \textquestiondown{}cuál es la probabilidad de que juegue al tenis?
\end{enumerate}
